\documentclass[UTF8,a4paper,11pt,openany]{ctexbook}
\usepackage{../common/statlearnbook}
\SLSetBookMetadata{第 2 册}{统计学习方法讲义 v2.6.0}{基础监督学习方法}
\begin{document}
\setcounter{page}{288}
\setcounter{chapter}{15}
\setcounter{figure}{0}
\pagestyle{headings}
\chaptermark{决策树}
\begin{optimizationbox}
每个结点执行有限候选集合上的局部最优化。局部最优切分不保证整棵树的全局最优，因此还需停止规则和剪枝控制模型复杂度。
\end{optimizationbox}
% v2.3.1 figure UID: FIG-P242-01
% v2.6.0 Image-reference reverse redraw: align tree paths with their axis-aligned regions.
\newcommand{\SLTreePartitionTitleStyle}{\fontsize{10.5pt}{12.6pt}\selectfont\bfseries}
\tikzset{
  slfig-FIG-P242-01/.style={font=\fontsize{9.2pt}{11.0pt}\selectfont,
    every path/.append style={line cap=round,line join=round}},
  slfig-FIG-P242-01-gold-edge/.style={draw=SLGold!82!black,line width=1.02pt,-{Stealth[length=1.9mm,width=1.2mm]}},
  slfig-FIG-P242-01-branch/.style={font=\fontsize{9.2pt}{11pt}\selectfont},
  slfig-FIG-P242-01-gold-branch/.style={slfig-FIG-P242-01-branch,text=SLGold!82!black},
  slfig-FIG-P242-01-x-split/.style={draw=SLBlue,line width=.82pt},
  slfig-FIG-P242-01-y-split/.style={draw=SLTeal,line width=.78pt,dash pattern=on 3.2pt off 2.0pt},
  slfig-FIG-P242-01-highlight/.style={draw=SLGold!82!black,line width=1.02pt},
  slfig-FIG-P242-01-region-label/.style={font=\fontsize{9.5pt}{11.4pt}\selectfont\bfseries}
}
\forestset{
  slfig-FIG-P242-01-tree/.style={
    for tree={font=\fontsize{9.2pt}{11pt}\selectfont,draw=SLBlue,fill=SLBlueBG,
      rounded corners=1.5mm,line width=.80pt,inner xsep=4pt,inner ysep=3pt,
      edge={draw=SLBlue,line width=.82pt,-{Stealth[length=1.9mm,width=1.2mm]}}}
  }
}
\pgfplotsset{slfig-FIG-P242-01-axis/.style={
  axis line style={draw=SLInk,line width=.55pt},tick style={draw=SLInk,line width=.52pt},
  tick label style={font=\fontsize{8.7pt}{10.4pt}\selectfont},
  label style={font=\fontsize{9.5pt}{11.4pt}\selectfont},clip=false}}
\begin{figure}[H]
\centering
\begin{minipage}[c]{.48\linewidth}
\centering
{\SLTreePartitionTitleStyle 决策树：条件与分支值}\par\smallskip
\begin{forest}
  slfig tree,
  slfig-FIG-P242-01-tree,
  for tree={font=\footnotesize,l sep=9mm,s sep=6mm,minimum width=17mm},
  [{$x_1\le2.4$}
    [{$R_1$},edge label={node[midway,left,slfig-FIG-P242-01-branch]{是}},slfig pattern blue]
    [{$x_2\le1.6$},edge={slfig-FIG-P242-01-gold-edge},
      edge label={node[midway,right,slfig-FIG-P242-01-gold-branch]{否}}
      [{$R_2$},slfig pattern gold,
        edge={slfig-FIG-P242-01-gold-edge},
        edge label={node[midway,left,slfig-FIG-P242-01-gold-branch]{是}}]
      [{$R_3$},slfig pattern teal,edge label={node[midway,right,slfig-FIG-P242-01-branch]{否}}]
    ]
  ]
\end{forest}
\par\smallskip
{\centering\fontsize{9.2pt}{11pt}\selectfont
\textcolor{SLGold!82!black}{金色路径：}$x_1>2.4$ 且 $x_2\le1.6$\par}
\end{minipage}\hfill
\begin{minipage}[c]{.48\linewidth}
\centering
{\SLTreePartitionTitleStyle 特征空间：同一轴对齐分裂}\par\smallskip
\begin{tikzpicture}[slfig-FIG-P242-01]
\begin{axis}[slfig axis,slfig-FIG-P242-01-axis,
  width=\linewidth,height=.84\linewidth,
  xmin=0,xmax=4,ymin=0,ymax=3,axis equal image,
  axis lines=box,xlabel={$x_1$},ylabel={$x_2$},
  xtick={2.4},ytick={1.6},clip=false,
]
  \path[fill=SLBlue!6,slfig pattern blue] (axis cs:0,0) rectangle (axis cs:2.4,3);
  \path[fill=SLGold!8,slfig pattern gold] (axis cs:2.4,0) rectangle (axis cs:4,1.6);
  \path[fill=SLTeal!6,slfig pattern teal] (axis cs:2.4,1.6) rectangle (axis cs:4,3);
  \addplot[slfig-FIG-P242-01-x-split] coordinates {(2.4,0) (2.4,3)};
  \addplot[slfig-FIG-P242-01-y-split] coordinates {(2.4,1.6) (4,1.6)};
  \draw[slfig-FIG-P242-01-highlight] (axis cs:2.4,0) rectangle (axis cs:4,1.6);
  \node[slfig-FIG-P242-01-region-label,text=SLBlue] at (axis cs:1.2,1.5) {$R_1$};
  \node[slfig-FIG-P242-01-region-label,text=SLGold!82!black] at (axis cs:3.2,.8) {$R_2$};
  \node[slfig-FIG-P242-01-region-label,text=SLTeal] at (axis cs:3.2,2.3) {$R_3$};
\end{axis}
\end{tikzpicture}
\end{minipage}
\caption{决策树中的每条根到叶路径对应特征空间中的一个轴对齐区域。}\label{fig:V2-C04-tree-partition}
\end{figure}

图\ref{fig:V2-C04-tree-partition}左侧的每条根到叶路径都对应右侧一个互斥区域；叶结点输出只在该区域内生效。
\begin{example}[八样本二叉划分的信息增益]
父结点有8个样本，正负类各4个。特征$A$把样本分成两个容量均为4的子集，求特征$A$的信息增益。
\end{example}
\end{document}
